\chapter{RELATED WORK}

 \section{Fine-Grained Visual Classification}

Fine-grained visual classification (FGVC) refers to the task of distinguishing subordinate categories that belong to the same general visual category, such as bird species, aircraft variants, or car models \cite{Fine-grainedSurvey}. Compared with generic object recognition, FGVC is more challenging because the visual differences between classes are often subtle and localized, while images from the same class may still vary significantly due to pose, viewpoint, illumination, background, and object deformation \cite{Fine-grainedSurvey, Zhou_2016_CVPR, TransIFC}. In bird species recognition, two species may share nearly identical global shapes and color distributions, differing only in minute visual traits such as bill shape, wing pattern, crown color, or breast pattern \cite{wah2011cub,HierarchicalFeatureAtt, He_2017}. Consequently, effective FGVC systems must be capable of learning representations that simultaneously capture global object structure and localized discriminative details \cite{HierarchicalFeatureAtt}.

Early FGVC studies often relied on part annotations, bounding boxes, or explicitly designed part localization mechanisms to compare corresponding regions across images. For example, part-based R-CNN methods learn object and part detectors and use geometric constraints to construct pose-normalized representations for fine-grained categorization \cite{zhang2014part}. More recent deep learning methods reduce the dependency on manual part supervision by learning discriminative regions or feature interactions directly from image-level labels. Representative approaches include bilinear feature modeling \cite{lin2015bilinear}, recurrent attention \cite{fu2017racnn}, and region proposal-based attention mechanisms \cite{yang2018learning}. These methods demonstrate that discovering local discriminative evidence is central to fine-grained recognition. However, many conventional FGVC methods are formulated as single-level fine-grained classification problems, where the output is only the finest category and higher-level taxonomic consistency is not explicitly modeled \cite{Fine-grainedSurvey, park2025visually}.

This thesis focuses on fine-grained bird recognition in a hierarchical setting. Instead of predicting only the species label, the model predicts labels at multiple taxonomic levels, namely order, family, and species. This formulation is particularly relevant for bird classification, where the taxonomy provides useful semantic structure and where mistakes between closely related species are common. The following subsections review the CUB-200-2011 dataset, discriminative region learning in FGVC, and diversity regularization methods that motivate the attribute-guided learning strategy used in this thesis.


\subsection{Bird Species Recognition and the CUB-200-2011 Dataset}

The Caltech-UCSD Birds-200-2011 dataset \cite{wah2011cub}, commonly known as CUB-200-2011, is one of the most widely used benchmarks for fine-grained bird image classification. It contains 200 bird species and 11,788 images, with annotations including bounding boxes, part locations, and 312 binary visual attributes. These attributes describe semantic visual properties of birds, such as color, shape, and pattern, and are organized around bird parts such as the bill, wing, breast, belly, tail, head, and leg.

CUB-200-2011 is particularly suitable for studying attribute-guided fine-grained classification for several reasons. First, the species are visually similar and often differ only in localized traits. Second, the dataset provides rich attribute annotations that can be used as auxiliary semantic supervision. Third, the bird taxonomy can be organized into a hierarchy of order, family, and species, making it appropriate for hierarchical image classification. In this thesis, the CUB-200-2011 hierarchy follows the H-CAST setting with 13 orders, 38 families, and 200 species \cite{park2025visually}.

The attributes in CUB are useful because they provide human-interpretable descriptions of subtle visual cues. For instance, two bird species may have a similar body shape and overall color, but differ in bill shape, wing color, crown pattern, or belly pattern. These attributes can guide a model toward fine-grained traits that may be difficult to learn from class labels alone, consistent with prior attribute-guided and attribute-aware attention methods for fine-grained recognition \cite{liu2017localizing, han2018attribute}. In this thesis, CUB attributes are not used as manual inputs during inference. Instead, they are used during training as auxiliary supervision for learning attribute-oriented representations. This design preserves the standard image-only inference setting while exploiting additional semantic information available in the training data.

Nevertheless, using all available attributes may not always be optimal. Some attributes may be redundant, weakly discriminative, or irrelevant for distinguishing a particular group of sibling species. This motivates the study of attribute selection strategies, including the proposed Hard-Sibling Discriminative Attribute Selection (HSDAS), which aims to identify attributes that are useful for distinguishing taxonomically related and visually similar bird classes.


\subsection{Discriminative Region Learning in FGVC}

A central challenge in FGVC is learning where to look. Since fine-grained categories are often separated by small local differences, many methods aim to identify discriminative object regions or part-level patterns. Bilinear CNN models, for example, model pairwise interactions between local convolutional features and have been shown to be effective for texture-rich and fine-grained recognition tasks \cite{lin2015bilinear}. Although bilinear pooling does not explicitly require part annotations, it can capture local feature interactions that are useful for distinguishing visually similar categories.

Another line of work uses attention mechanisms to locate informative regions. RA-CNN recursively learns region attention and region-based feature representations at multiple scales, allowing the model to progressively focus on discriminative image regions \cite{fu2017racnn}. Similarly, NTS-Net introduces a Navigator-Teacher-Scrutinizer framework, where informative regions are proposed and then scrutinized for final fine-grained prediction \cite{yang2018learning}. These methods highlight the importance of localized evidence for fine-grained recognition.

Despite their effectiveness, many discriminative region learning methods are designed for flat species-level classification. They aim to improve the prediction of the finest label, but they do not explicitly model taxonomy-level consistency. In contrast, hierarchical classification requires the model to produce predictions at multiple semantic levels. For example, a species prediction should be compatible with its corresponding family and order. This thesis therefore builds on the visually consistent hierarchical classification framework of H-CAST, which uses segment-based visual grouping to align fine-to-coarse predictions \cite{park2025visually}.

The proposed approach differs from conventional discriminative region learning in two ways. First, it uses hierarchical visual representations from H-CAST rather than a flat image classifier. Second, it incorporates bird attributes as auxiliary supervision, encouraging the model to learn representations related to semantic visual traits. Instead of directly predicting attributes as the final task, the model uses attribute-oriented features to support species-level classification while preserving hierarchical predictions at the family and order levels.


\subsection{Loss-based Diversity Regularization}

In addition to architecture-based region learning, another direction in FGVC encourages discriminative and diverse feature learning through loss functions. Mutual-Channel Loss (MC-Loss) is a representative example of this direction \cite{Chang_2020}. MC-Loss was proposed to learn localized discriminative features without requiring explicit part or bounding-box annotations. It consists of two main components: a discriminality component and a diversity component. The discriminality component encourages feature channels associated with a class to be individually discriminative, while the diversity component encourages those channels to focus on different spatial regions.

The intuition behind the diversity component is that a fine-grained object cannot usually be recognized from a single local cue. For bird recognition, different parts such as the head, wing, breast, bill, and tail may all provide complementary evidence. If all feature channels attend to the same most salient region, the model may ignore other useful visual traits. By encouraging feature channels to be spatially diverse, MC-Loss helps the network discover multiple discriminative regions and reduce redundant feature responses \cite{Chang_2020}.

This idea is relevant to attribute-guided bird classification. In the proposed model, different attribute groups may correspond to different semantic parts or visual properties. For example, bill shape, crown color, wing pattern, and tail shape should ideally provide complementary information rather than collapsing onto the same image region. Therefore, this thesis adopts the diversity principle from MC-Loss and adapts it to the attribute-guided H-CAST setting.

It is important to note that the diversity regularization used in this thesis is not identical to the original MC-Loss. MC-Loss applies diversity regularization to class-specific convolutional feature channels. In contrast, this thesis applies a diversity-inspired regularization term to attribute-group attention patterns. The goal is to encourage different attribute groups or macro-parts to attend to complementary visual regions. Thus, the proposed loss should be understood as an adaptation of the MC-Loss diversity idea, rather than a direct reuse of the original MC-Loss formulation.

By combining attribute supervision with diversity regularization, the model is encouraged to learn attribute-oriented representations that are both discriminative and complementary. This is especially important for CUB-200-2011, where species-level differences often depend on multiple subtle visual attributes. The effect of this diversity regularization is evaluated in later experimental chapters through ablation studies with different loss weights.

 \section{Attribute-Based and Prompt-Assisted Fine-Grained Recognition}

Visual attributes provide a semantic description of object appearance. Instead of representing an object only by its category label, attributes describe human-interpretable properties such as color, shape, texture, pattern, or part configuration. In image recognition, this representation is useful because many visual categories share common properties. For example, different bird species may share attributes such as a pointed bill or striped wings, while still differing in other subtle traits. Therefore, attributes can serve as an intermediate semantic layer between low-level visual features and high-level class labels.

Attribute-based recognition has been studied as a way to make visual models more interpretable and transferable. Farhadi et al. introduced the idea of describing objects by attributes, shifting recognition from only naming object categories to also describing their visual properties \cite{farhadi2009attributes}. Lampert et al. further showed that attributes can support recognition of unseen categories by transferring knowledge through semantic attribute descriptions \cite{lampert2009attributetransfer, lampert2014attributebased}. Although these works are often discussed in the context of zero-shot learning, the underlying idea is also relevant to supervised fine-grained recognition: attributes can guide a model toward visual traits that are shared across classes and meaningful to humans.

In this thesis, attributes are not used for zero-shot prediction. Instead, they are used as auxiliary supervision for learning better image representations. The goal is to encourage the model to encode fine-grained visual traits that are useful for distinguishing bird species. This differs from attribute-based zero-shot methods that rely on class-level attribute prototypes or semantic descriptions as part of the inference mechanism. In the proposed framework, no manually provided image-level attribute labels are required at inference time.


\subsection{Visual Attributes as Semantic Supervision}

In conventional supervised image classification, the training signal is usually provided only by class labels. This supervision tells the model which category an image belongs to, but it does not explicitly describe why the image belongs to that category. For fine-grained classification, this limitation is important because two classes may be separated by very small differences. A class label alone may not explicitly encourage the model to focus on the bird's bill, crown, wing, tail, or breast pattern. Visual attributes provide additional semantic supervision by describing these discriminative properties.

Attributes can be used in several different ways. One common strategy is to predict attributes as intermediate semantic outputs and then use the predicted attributes for classification. The Direct Attribute Prediction and Indirect Attribute Prediction frameworks of Lampert et al. are representative examples in zero-shot recognition \cite{lampert2009attributetransfer, lampert2014attributebased}. Another strategy is to use attributes as class-level semantic embeddings. Akata et al. embed class labels in an attribute-vector space and learn image-to-label compatibility for attribute-based classification \cite{akata2013labelembedding, Akata_2016}. Semantic Autoencoder later followed the semantic embedding direction by learning a projection from visual features to semantic space and adding a reconstruction constraint from semantic space back to visual features \cite{kodirov2017semantic}. These methods demonstrate that attributes can define a structured semantic space for visual recognition.

Attributes can also be used as auxiliary supervision or as guidance for attention. Liu et al. proposed an attribute-guided attention localization method for fine-grained recognition, where part attribute descriptions guide the learning of local region localizers on CUB-200-2011 \cite{liu2017localizing}. Their method trains localizers for parts such as head, wing, breast, and tail using part attribute descriptions, and the attribute descriptions are used during training rather than as test-time inputs. Han et al. proposed the Attribute-Aware Attention Model (A$^3$M), which learns local attribute representations and global category representations simultaneously through attribute-guided and category-guided attention modules \cite{han2018attribute}. Huynh and Elhamifar proposed dense attribute-based attention for fine-grained generalized zero-shot learning, where each attribute is associated with relevant image regions to obtain attribute-based features \cite{huynh2020DenseAttribute}. These works provide concrete examples of how attributes can guide the model toward local discriminative regions.

The auxiliary-supervision setting is particularly relevant to this thesis. Instead of using attributes as the final prediction target or requiring human-provided attributes at inference time, the proposed model uses attribute labels during training to shape the learned representation. This design is related to previous attribute-guided attention methods, but differs in two aspects. First, it is built on an H-CAST-based hierarchical visual backbone rather than a flat image classifier. Second, it uses attribute-group queries and species-level attribute fusion to support hierarchical bird classification.

The main advantage of attribute supervision is that it introduces structured semantic information into the learning process. Instead of treating each species as an unrelated index, the model can receive additional signals about visual properties that may be shared across species. For example, attributes related to wing color, bill shape, or belly pattern may help the model learn features that are useful for multiple bird classes. This can regularize the learned representation and make it more aligned with human-interpretable visual traits.


However, attribute supervision also introduces challenges. First, not all attributes are equally useful for a specific classification problem. Some attributes may be common across many classes and therefore weakly discriminative. Second, different attributes may be redundant, especially when they describe correlated visual properties. Third, noisy or irrelevant attributes may distract the model from the main classification objective. This concern is consistent with recent work on attribute-based visual recognition showing that very large automatically generated attribute sets can contain substantial noise, and that concise attribute subsets may preserve performance while improving interpretability \cite{yan2023learning}. These challenges motivate the attribute selection component studied later in this thesis.



\subsection{Prompt-Assisted Attribute Reasoning}

Recent prompt-based and language-assisted fine-grained recognition methods suggest that semantic visual descriptions can help when direct subordinate-level recognition is difficult. For example, FineR uses visual question answering models to extract part-level visual attributes, such as wing color, head pattern, and tail pattern, and then uses large language models to reason about possible fine-grained category names from these semantic descriptions \cite{liu2024democratizing}. Related work on fine-grained visual descriptions also shows that textual or semantic descriptions can provide useful side information for representing fine-grained categories \cite{reed2016learning}. These studies support the view that attributes can act as intermediate semantic evidence between raw image features and fine-grained class decisions.

However, this thesis studies a different setting. It does not use prompts, large language models, or manually provided attribute descriptions at inference time. Instead, it uses the structured CUB visual attributes as training-time auxiliary supervision in a supervised hierarchical classification framework \cite{wah2011cub}. Thus, prompt-assisted recognition motivates the usefulness of attributes, while the proposed method studies how such information can be incorporated into an image-only hierarchical classifier.



\subsection{Attributes in Fine-Grained Bird Recognition}

CUB-200-2011 is especially suitable for studying attributes in fine-grained bird recognition because it provides both species labels and rich visual attribute annotations \cite{wah2011cub}. The dataset contains 200 bird species and 312 binary attributes describing local and global bird appearance. These attributes include properties related to bill shape, wing color, upperparts color, underparts color, breast pattern, tail shape, head pattern, eye color, leg color, and other bird parts. Such annotations make CUB a natural benchmark for attribute-guided fine-grained classification. In this thesis, these annotations are used through the class-level continuous attribute file, where each value represents the percentage of annotators who considered an attribute present for a given species.

Bird species recognition benefits from attributes because many species are distinguished by localized traits. For instance, two species may have similar body shapes and overall color distributions, but differ in crown color, wing pattern, bill length, or tail shape. These differences are often subtle and may occupy only a small part of the image. Attribute supervision can therefore provide a useful learning signal by indicating which kinds of visual traits are relevant for discrimination.

Prior work has also shown the importance of semantic descriptions for fine-grained visual recognition. Reed et al. studied fine-grained visual descriptions and learned joint representations between images and textual descriptions on datasets including CUB \cite{reed2016learning}. Their work demonstrates that semantic side information can be useful for representing fine-grained categories. In contrast to description-based or zero-shot recognition settings, this thesis uses the structured attribute annotations already provided by CUB as auxiliary supervision for a supervised hierarchical classification task.

The use of CUB attributes in this thesis differs from direct attribute-based classification. The model does not first predict a complete attribute vector and then classify the bird solely from that vector. Instead, attribute-related supervision is integrated into an H-CAST-based visual backbone. Attribute-group queries are used to extract attribute-oriented features from species-level graph-node representations, and these features are fused with the species-level class token. The attribute predictions act as auxiliary outputs during training, while the final inference remains image-based.

Another important issue is attribute subset selection. Since CUB contains 312 attributes, using all attributes may introduce redundancy or weakly useful supervision. Some attributes may help distinguish only a few species, while others may be common across many sibling classes. This thesis therefore compares multiple attribute usage strategies, including using all attributes, an internal node-vs-siblings statistical selection baseline, random attribute subsets, and the proposed Hard-Sibling Discriminative Attribute Selection (HSDAS). The internal statistical baseline, designed in this thesis and referred to as Sibling-IVAS-188, uses train-only image-level binary attributes and ranks attributes for each taxonomy node primarily by positive occurrence-rate difference, with Information Value used as a secondary criterion. The goal is to study whether a compact and taxonomically meaningful attribute subset can better support hierarchical fine-grained bird classification.


 \section{Hierarchical Image Classification}

Hierarchical image classification \cite{silla2011survey,chang2021your,jiang2024hierarchical} extends standard image classification by predicting labels according to a predefined class hierarchy. Instead of assigning each image to only one flat class, a hierarchical classifier can produce predictions at multiple semantic levels, such as order, family, and species in bird recognition. This setting is useful when the label space naturally follows a taxonomy and when different levels of granularity are meaningful for different users or application scenarios.


In the context of this thesis, hierarchical classification is particularly relevant because bird categories are naturally organized taxonomically. A model that predicts only the species label may achieve fine-grained recognition, but it does not explicitly evaluate whether the predicted species is compatible with the corresponding family and order. Therefore, hierarchical classification provides a more structured formulation for bird recognition, where both fine-grained accuracy and taxonomy-level consistency are important.

\subsection{Flat Classification vs. Hierarchical Classification}

Flat classification treats all target categories as independent classes at a single semantic level \cite{silla2011survey}. In fine-grained bird classification, this usually means predicting one of the 200 species in CUB-200-2011. This formulation is simple and widely used, but it ignores the hierarchical relationships between classes. For example, if a model predicts a species from one family while its family-level prediction indicates a different family, the prediction path becomes inconsistent with the taxonomy.

Hierarchical classification, by contrast, explicitly models labels at multiple levels of abstraction. In a bird taxonomy, the model may predict the order, family, and species of an image. This formulation allows the model to represent both coarse and fine semantic information. It also makes evaluation more informative, since hierarchy-aware evaluation treats an error between two species in the same family as less severe than an error between species from different orders \cite{bertinetto2020making}.

Existing hierarchical classification methods can be broadly divided into single-level and full-taxonomy prediction approaches. In a bottom-up single-level approach, the model predicts the finest label, and higher-level labels are inferred from the taxonomy. This guarantees a valid taxonomy path, but it relies heavily on the correctness of the fine-grained prediction. In a top-down approach, the model first predicts a coarse label and then refines the prediction toward finer labels. This can be useful when coarse-level recognition is easier, but early mistakes may propagate downward \cite{park2025visually}.

Full-taxonomy prediction aims to predict all hierarchy levels simultaneously. This is the setting used in this thesis. It provides explicit supervision at each level and allows the model to evaluate performance at order, family, and species levels separately. However, because different heads may predict different levels independently, full-taxonomy prediction must address the problem of taxonomy consistency \cite{park2025visually}.

\subsection{Taxonomy Consistency and Multi-Granularity Prediction}

A central challenge in hierarchical classification is maintaining consistency among predictions at different levels. A prediction is taxonomically consistent if the predicted labels form a valid path in the hierarchy. For example, a predicted species should belong to the predicted family, and that family should belong to the predicted order. If the species and family predictions do not match the known taxonomy, the model produces an inconsistent prediction path.

This issue is especially important in multi-granularity fine-grained recognition. Different levels of the hierarchy may require different visual evidence. Coarse-level labels often depend on broad object-level cues, while fine-grained species labels may depend on subtle local details. H-CAST emphasizes this problem by showing that classifiers at different hierarchy levels may attend to different image regions, leading to inconsistent predictions even when semantic supervision is used \cite{park2025visually}.

Several metrics are used to evaluate hierarchical prediction. Level-wise accuracy measures the classification accuracy at each individual level, such as order accuracy, family accuracy, and species accuracy. However, level-wise accuracy alone does not fully capture whether the predictions form a valid taxonomy path. Tree-based InConsistency Error rate (TICE) measures the proportion of samples whose predicted path is inconsistent with the hierarchy. Full-Path Accuracy (FPA), on the other hand, counts a prediction as correct only when all hierarchy levels are predicted correctly. Therefore, FPA evaluates both correctness and hierarchical agreement, while TICE focuses on the validity of the predicted path \cite{park2025visually}.

Multi-granularity prediction is useful because it allows a model to provide predictions at different levels of specificity. In practical bird recognition, an image may not always contain enough detail for confident species-level prediction, but a coarse-level family or order prediction may still be meaningful. At the same time, predicting multiple levels introduces a trade-off: the model must learn features that are useful for coarse categories while preserving subtle visual cues needed for fine-grained species discrimination.

\subsection{Representative Hierarchical Classification Methods}

A number of methods have been proposed to exploit hierarchy information in image classification. Bertinetto et al. studied how class hierarchies can be used to make ``better mistakes'' by encouraging errors that are closer to the ground-truth class in the hierarchy \cite{bertinetto2020making}. Their work highlights that not all classification errors are equally severe and motivates the use of hierarchy-aware evaluation and training objectives.

For fine-grained visual classification, Chang et al. proposed FGN, which reformulates traditional FGVC as a multi-granularity recognition problem \cite{chang2021your}. Instead of predicting only the finest category, FGN considers a coarse-to-fine label hierarchy and uses level-specific classification heads to reduce negative transfer across granularity levels. The method also allows fine-grained features to contribute to coarse-level prediction, reflecting the observation that fine-level visual details can help recognize coarser labels.

Chen et al. proposed the Hierarchical Residual Network (HRN) for hierarchical multi-granularity classification \cite{Chen_2022_CVPR}. HRN models hierarchical feature interaction by adding granularity-specific features from parent levels to child-level features through residual connections. It also uses a combinatorial loss based on the tree hierarchy to exploit label relationships and transfer hierarchical knowledge during training \cite{Chen_2022_CVPR}. This approach focuses on transferring hierarchical knowledge between levels and improving feature interaction across granularities.

With the rise of Vision Transformers, Wang et al. introduced TransHP, a hierarchical prompting method for image classification \cite{wang2023transhp}. TransHP learns prompt tokens representing coarse ancestor classes and injects the predicted coarse-class prompt into intermediate transformer features. The motivation is that ancestor-level information can guide the model to focus on subtle differences among descendant classes. Unlike residual CNN-based approaches, TransHP uses the transformer token mechanism to condition feature extraction on hierarchy information.

Another related direction focuses directly on prediction consistency. Wang et al. proposed Consistency-Aware Feature Learning (CAFL) for hierarchical fine-grained visual classification \cite{wang2023consistency}. CAFL encodes prediction consistency constraints through forward deduction and backward induction over the label hierarchy, aiming to improve both classification accuracy and consistency. This type of method shows that enforcing consistency is important, but it mainly focuses on semantic relations between labels and prediction scores.

More recently, H-CAST introduced a visually consistent approach to hierarchical image classification \cite{park2025visually}. Unlike methods that rely mainly on label relations or semantic losses, H-CAST builds on CAST and uses superpixel-based hierarchical visual grouping. It applies fine-to-coarse supervision at different stages of the model so that finer visual segments support fine-grained labels and coarser segments support higher-level labels. In addition, H-CAST uses Tree-Path KL Divergence to encourage semantic consistency across hierarchy levels.

H-CAST is the most directly related baseline for this thesis. It provides a strong image-only framework for hierarchical bird classification and addresses both visual and semantic consistency. However, in its CUB setting, H-CAST uses the order-family-species taxonomy but does not incorporate the CUB visual attribute annotations as explicit auxiliary supervision. This thesis builds on the H-CAST framework by introducing attribute-guided learning at the species-level branch. The goal is to investigate whether semantic attribute supervision and hard-sibling-aware attribute selection can further improve fine-grained species discrimination while preserving hierarchical prediction.

 \section{Vision Transformers and Segment-Based Visual Grouping}

Vision Transformers have become an important backbone family for image recognition. Unlike convolutional neural networks, which process images through local convolutional filters and spatial downsampling, Vision Transformers represent an image as a sequence of tokens and use self-attention to model relationships among them \cite{dosovitskiy2020image}. This token-based formulation provides a flexible framework for visual representation learning, but standard ViT models usually rely on fixed square patch tokens. Such tokens are convenient for computation, but because they are defined by a regular grid, they may cut across object boundaries and are not explicitly designed to correspond to semantic parts \cite{ke2024learning,park2025visually}. 

For fine-grained and hierarchical recognition, this limitation is important. Distinguishing bird species often requires localized evidence from object parts, such as the bill, wing, head, breast, or tail. At the same time, higher-level categories such as family or order may depend on broader visual structures. Therefore, a model that can organize visual information from small parts to larger regions is naturally suitable for hierarchical bird classification. This motivates the use of segment-based visual grouping methods in CAST \cite{ke2024learning} and H-CAST \cite{park2025visually}.


\subsection{Vision Transformers}

The Vision Transformer (ViT) adapts the Transformer architecture \cite{vaswani2017attention} to image classification by dividing an image into fixed-size patches and treating each patch as a token \cite{dosovitskiy2020image}. A learnable class token is appended to the patch tokens, and the sequence is processed by multiple Transformer encoder blocks. Through self-attention, each token can interact with other tokens, allowing the model to capture long-range dependencies across the image.

ViT provides a simple and powerful image classification framework, but its patch-based representation has limitations for visual parsing. Since patches are defined by a regular grid, they may cut across object boundaries or merge foreground and background pixels. This can be problematic in fine-grained recognition, where the relevant visual evidence may correspond to irregularly shaped parts rather than square image regions.

Another limitation is that standard ViT maintains patch tokens as the main visual units throughout the network. Although self-attention can aggregate information globally, the model does not explicitly produce a part-to-whole segmentation hierarchy. For hierarchical classification, this means that the model may not directly align fine-grained visual cues with fine labels and broader visual regions with coarse labels. These issues motivate later segment-based extensions of ViT, including CAST and H-CAST.


\subsection{CAST}

The model is named CAST because it Concurrently learns segmentation and recognition using Adaptive Segment Tokens, modifying the standard ViT formulation by replacing fixed square patches with adaptive segment tokens based on superpixels \cite{ke2024learning}. CAST enhances the standard ViT formulation with adaptive segment tokens and graph pooling, allowing segmentation and recognition to be learned concurrently within a single model \cite{ke2024learning}. Its central idea is that segmentation should not be treated only as an independent end task, but as an internal visual parsing process that evolves with and supports image recognition. CAST learns hierarchical segmentation together with recognition using image-level recognition objectives, without requiring pixel-level segmentation labels.

CAST starts by extracting superpixels from the input image using SEEDS \cite{van2015seeds}. These superpixels serve as the finest visual units and are used to construct adaptive segment tokens. Compared with fixed patch tokens, superpixels can better follow low-level color and boundary structure than fixed square patches, providing adaptive visual units before higher-level graph pooling. Each segment token is then processed by Transformer blocks, and graph pooling is used to progressively merge finer segments into coarser segments.

The graph pooling module constructs a fine-to-coarse hierarchy of segment tokens. At each level, finer segment tokens are softly assigned to coarser segment tokens based on feature similarity. This produces a hierarchy of segmentations and features, allowing the model to gradually move from local parts to larger object regions. The CAST paper emphasizes that this hierarchical segmentation is not trained with external segmentation labels; instead, it emerges from the recognition objective itself \cite{ke2024learning}.

This design distinguishes CAST from standard ViT in two major ways. First, the visual tokens are adaptive segments rather than fixed patches. Second, the number and granularity of tokens change through graph pooling, producing a part-to-whole visual hierarchy. As a result, CAST can provide both image-level recognition and internal hierarchical segmentation. The CAST paper reports that this segment-based representation improves unsupervised part-whole discovery, semantic segmentation, image classification, and efficiency compared with ViT baselines \cite{ke2024learning}.

CAST is relevant to this thesis because it provides the visual backbone on which H-CAST is built. Its fine-to-coarse segment hierarchy makes it possible to connect visual granularity with semantic granularity. In bird classification, fine segments can preserve local details such as the head, bill, or wing, while coarser segments can represent larger bird regions. This property is useful for hierarchical classification, where different taxonomy levels may require different levels of visual detail.


\subsection{H-CAST}

H-CAST extends CAST to hierarchical image classification \cite{park2025visually}. While CAST learns hierarchical segmentation for a flat recognition objective, H-CAST introduces hierarchical supervision so that different stages of the visual hierarchy are associated with different taxonomy levels. The model is motivated by the observation that fine-grained and coarse-grained classification may rely on different but related visual cues. For example, distinguishing between visually similar bird species may require local details such as the bill, wings, or head coloration, whereas recognizing a higher-level category may require broader object-level evidence.

The key idea of H-CAST is internal visual consistency. H-CAST argues that predictions at different hierarchy levels should be grounded in coherent and related visual regions. H-CAST addresses this issue by using the fine-to-coarse visual grouping structure inherited from CAST. Fine-level visual details are progressively grouped into coarser regions, so predictions at different hierarchy levels are encouraged to rely on coherent and related image regions.

Formally, H-CAST considers an image associated with labels at multiple hierarchy levels. It uses the segment and feature representations at different levels of the CAST hierarchy, adds a classification head at each level, and applies cross-entropy supervision across all hierarchy levels. In contrast to CAST, which uses a flat supervision signal on the final class token, H-CAST applies hierarchical supervision at multiple stages of the visual parsing process \cite{park2025visually}. This allows fine and coarse labels to guide different levels of visual representation.

In addition to visual consistency, H-CAST introduces Tree-Path KL Divergence loss to improve semantic consistency. This loss concatenates the labels from all hierarchy levels into a single distribution and aligns the concatenated outputs of the classification heads with this distribution using KL divergence. The purpose is to reduce prediction paths that violate the taxonomy, such as predicting a species whose parent family or order is inconsistent with the predicted higher-level labels \cite{park2025visually}.

The H-CAST paper evaluates hierarchical classification using both accuracy and consistency metrics. It defines Full-Path Accuracy (FPA) as the proportion of samples for which all hierarchy levels are correctly predicted, and uses Tree-based InConsistency Error rate (TICE) to measure whether predicted paths are valid in the taxonomy \cite{park2025visually}. These metrics are also relevant to this thesis because the proposed model is evaluated not only by species accuracy, but also by family accuracy, order accuracy, FPA, and TICE.

H-CAST is the main baseline and architectural foundation of this thesis. The proposed method keeps the H-CAST visual parsing pipeline and extends it with attribute-guided learning for CUB-200-2011. Specifically, instead of changing the hierarchical segment grouping mechanism, this thesis adds an attribute-guided module at the species-level branch. The motivation is that H-CAST provides visually consistent hierarchical representations, while CUB attributes provide additional semantic supervision about fine-grained bird traits. Combining these two sources of information may help the model better distinguish visually similar sibling species.

 \section{Attribute Selection for Fine-Grained Hierarchical Classification}

Attribute annotations provide useful semantic supervision for fine-grained recognition, but using all available attributes is not always optimal, since feature-selection studies emphasize the importance of removing irrelevant or redundant variables \cite{peng2005feature}, and recent work on attribute-based visual recognition also suggests that concise attribute subsets can preserve performance while improving interpretability \cite{yan2023learning}. In datasets such as CUB-200-2011, the attribute space is rich and detailed, covering color, shape, pattern, and body-part-related properties \cite{wah2011cub}. However, not every attribute is equally useful for distinguishing every bird species. Some attributes may be shared by many species, some may be redundant with other attributes, and some may be weakly related to the most difficult class distinctions. Therefore, when attributes are used as auxiliary supervision, attribute selection becomes an important design choice.

In general machine learning, feature selection aims to identify a subset of informative variables while removing irrelevant or redundant ones. Classical approaches often evaluate features according to their relevance to target labels and, in some cases, their redundancy with other features. For example, the minimum Redundancy Maximum Relevance (mRMR) criterion selects features that are highly relevant to the target while being minimally redundant with already selected features \cite{peng2005feature}. Relief-based methods estimate the quality of attributes by comparing nearby samples from the same and different classes, making them sensitive to local class distinctions \cite{kira1992practical, kononenko1994estimating}.In applied scorecard and risk modeling, Weight of Evidence and Information Value are commonly used during variable binning and screening to assess the usefulness of predictive variables \cite{siddiqi2012credit}.


Although these feature selection methods provide useful general principles, fine-grained hierarchical classification introduces additional considerations. First, the most useful attributes may not be globally discriminative across all classes. An attribute that separates visually unrelated birds may be less useful for distinguishing two sibling species within the same family. Second, the hierarchical structure of the labels means that discrimination can occur at multiple levels, such as order, family, and species. Third, in fine-grained bird classification, difficult errors often occur between visually and taxonomically related classes. Therefore, an effective attribute subset should support local discrimination among related classes, rather than only global separation across the entire dataset.


This idea is related to prior work on taxonomy-aware attribute querying. Mall et al. proposed a field-guide-inspired zero-shot learning framework in which a learner asks for the most useful attributes to describe a new class \cite{mall2021field}. Their method is motivated by real field guides, where a species can be described by a small number of informative attributes rather than requiring annotation of the complete attribute vector. In particular, their \emph{sibling-variance} strategy selects attributes that vary within a branch of the taxonomy, showing that taxonomy-aware attribute selection can be more informative than measuring attribute variation globally. This work is closely related to the motivation of this thesis, because both approaches emphasize attributes that distinguish related classes.


More broadly, attribute selection in this thesis is also related to hierarchical feature selection. Zhao et al. proposed hierarchical feature selection with recursive regularization, where different feature subsets are selected for different nodes in a class hierarchy by using parent-child and sibling relationships \cite{zhao2017hierarchical}. This line of work shows that hierarchical class structure can be useful for selecting features, because related nodes may share common features while sibling nodes may require discriminative features that separate them. However, these methods are general feature selection frameworks and are not specifically designed for CUB bird attributes or for attribute-guided visual representation learning.


This thesis considers attribute selection in the context of attribute-guided H-CAST. Since the proposed model uses attribute groups as auxiliary supervision and species-level feature guidance, the selected attributes influence which semantic cues are emphasized during training. If too many weak or redundant attributes are included, the attribute branch may learn less useful supervision and reduce the effectiveness of attribute-guided fusion. Conversely, if too few attributes are selected, the model may lose important cues needed to distinguish similar bird species. Attribute selection is therefore studied as a trade-off between compactness, discriminative power, and semantic coverage.

Several attribute usage strategies are considered in this thesis. The simplest strategy, referred to as \textbf{FullAttr-312}, uses all 312 CUB attributes. This setting preserves the full semantic information available in the dataset, but it may also include redundant or weakly discriminative attributes.

This thesis also includes an internal statistical attribute-selection baseline, referred to as \textbf{Sibling-IVAS-188}. This baseline was implemented as an exploratory comparison method rather than adopted from a published work. It uses train-only image-level binary attribute annotations and compares each taxonomy node with a corresponding negative set. For species-level nodes, the negative set consists of sibling species within the same family; for family-level nodes, it consists of sibling families within the same order; and for order-level nodes, it consists of the remaining orders. Attributes are ranked primarily by positive occurrence-rate difference between the target node and the negative set, with Information Value used as an additional ranking criterion. Therefore, Sibling-IVAS-188 should be understood as a taxonomy-aware statistical baseline, not as a global attribute-selection method.



To provide a stronger control, this thesis also includes \textbf{Random-199}, a random attribute subset with the same number of attributes as the main selected subset. This baseline is important because it helps determine whether any performance change is caused by the attribute selection criterion itself or merely by reducing the number of attributes. If a selected subset outperforms a random subset of the same size, this provides evidence that the selection strategy captures meaningful discriminative information.

The main attribute selection strategy studied in this thesis is \textbf{Hard-Sibling Discriminative Attribute Selection (HSDAS)}. HSDAS follows the high-level intuition of sibling-aware attribute selection: attributes should be chosen according to local taxonomic distinctions, not only according to global discriminative strength. However, the setting is different from field-guide-inspired zero-shot learning. Mall et al. use sibling-aware attribute querying to reduce the cost of expert annotation for novel classes \cite{mall2021field}. In contrast, this thesis assumes that class-level CUB attributes are already available during training and uses HSDAS to construct an attribute supervision mask for an H-CAST-based hierarchical classifier.

Conceptually, HSDAS reflects the nature of fine-grained bird classification. Two species from the same family may share global shape and overall color but differ in localized properties such as crown color, bill shape, wing pattern, or tail shape. Attributes that distinguish these hard sibling classes are more likely to support species-level recognition than attributes that separate visually unrelated birds. Therefore, HSDAS prioritizes attributes that distinguish a taxonomic node from its sibling nodes, especially hard siblings with similar attribute profiles.

In this thesis, HSDAS is not presented as a general-purpose feature selection algorithm for all datasets. Rather, it is designed for the specific setting where a taxonomy and class-level attribute annotations are available. Its role is to select a compact set of semantic attributes that can guide an H-CAST-based classifier toward discriminative bird traits. The selected attributes are then used to define active attribute groups for auxiliary attribute prediction and attribute-guided species-level fusion.

The study of attribute selection also supports the interpretability of the proposed approach. Since selected attributes correspond to human-readable bird traits, the resulting subset can be analyzed in terms of body parts and semantic groups. This makes it possible to inspect whether the model is guided by meaningful attributes such as bill shape, wing color, tail pattern, or breast markings. However, attribute selection does not imply that the model receives manually provided attribute labels at inference time. The attributes are used only during training as auxiliary supervision, while inference remains image-based.

Overall, attribute selection is a key component of the proposed attribute-guided hierarchical classification framework. It connects the semantic structure of CUB attributes with the taxonomic structure of bird labels. By comparing FullAttr-312, Sibling-IVAS-188, Random-199, and HSDAS-199, this thesis evaluates whether a compact and hard-sibling-aware attribute subset can better support fine-grained species recognition and hierarchical prediction than using all attributes, random subsets, or earlier internal attribute-selection baselines.


\section{Research Gap}

The previous sections show that fine-grained bird classification, attribute-based recognition, prompt-assisted fine-grained reasoning, and hierarchical image classification have each been studied from different perspectives. Fine-grained visual classification methods focus on subtle visual differences between closely related categories \cite{Fine-grainedSurvey, Zhou_2016_CVPR, TransIFC}. Attribute-based recognition methods show that semantic visual properties can describe object appearance and provide additional supervision or transferable representations \cite{farhadi2009attributes, lampert2014attributebased}. Prompt-assisted fine-grained recognition further suggests that semantic visual descriptions can serve as intermediate evidence when direct subordinate-level recognition is difficult \cite{liu2024democratizing}. Hierarchical classification methods exploit taxonomy structures to produce predictions at multiple semantic levels \cite{silla2011survey, chang2021your, Chen_2022_CVPR, wang2023transhp, park2025visually}. However, the integration of these directions remains limited in supervised fine-grained hierarchical bird classification.

First, image-level class supervision alone may not explicitly encourage a model to encode the semantic traits that distinguish visually similar species. This is important in CUB-200-2011, where many bird species belong to the same family and share highly similar global appearance. Their differences may appear only in subtle visual traits such as bill shape, wing pattern, head pattern, breast color, tail shape, or crown color \cite{wah2011cub}. Prior attribute-based and attribute-guided recognition studies show that attributes can provide semantic descriptions, transferable representations, or attention guidance for fine-grained recognition \cite{farhadi2009attributes, lampert2014attributebased, liu2017localizing, han2018attribute, huynh2020DenseAttribute}. Nevertheless, these methods are usually designed for zero-shot recognition, flat fine-grained classification, or attribute-aware attention, rather than supervised hierarchical bird classification with order, family, and species predictions.

Second, recent prompt-based and language-assisted fine-grained recognition methods highlight the value of semantic visual descriptions for fine-grained reasoning \cite{liu2024democratizing}. These methods suggest that when direct category prediction is difficult, intermediate descriptions of visual traits can help narrow down possible subordinate categories. However, such methods usually rely on prompts, large language models, or generated textual descriptions. In contrast, this thesis does not use prompts or manually provided attribute descriptions at inference time. The gap addressed here is how to use the structured CUB visual attributes as training-time auxiliary supervision so that the final model remains image-only during inference.

Third, visually consistent hierarchical models such as H-CAST provide strong image-based backbones for hierarchical prediction. H-CAST builds on CAST and uses fine-to-coarse segment grouping to align visual granularity with taxonomy levels, while Tree-Path KL Divergence encourages semantic consistency across predicted hierarchy levels \cite{ke2024learning, park2025visually}. However, in its CUB setting, H-CAST uses the order-family-species taxonomy but does not explicitly incorporate CUB visual attributes as auxiliary semantic supervision. This motivates the study of attribute-guided learning on top of an H-CAST-based hierarchical backbone, especially at the species-level branch where fine-grained visual traits are most important.

Fourth, even when attributes are available, it is unclear whether using all attributes is the best strategy for supervised hierarchical bird classification. CUB contains 312 attributes \cite{wah2011cub}, but not all of them are equally useful for every classification decision. Some attributes may be redundant, broadly shared across many species, or weakly related to the most difficult sibling-species confusions. General feature-selection methods provide criteria for selecting informative variables and reducing redundancy \cite{peng2005feature, kira1992practical, kononenko1994estimating}, but they do not directly address how human-interpretable bird attributes should be selected as auxiliary supervision inside a hierarchical visual classification model.

More closely related to this thesis, sibling-aware attribute querying has been studied in field-guide-inspired zero-shot learning. Mall et al. propose sibling-aware attribute querying strategies that prioritize attributes useful for distinguishing a class from related classes in a taxonomy, showing that local taxonomic distinctions can be more informative than global attribute variation \cite{mall2021field}. However, their setting focuses on reducing expert annotation cost for zero-shot class descriptions. In contrast, this thesis studies a supervised hierarchical classification setting where class-level CUB attributes are already available during training. The gap addressed here is how to use hard-sibling-aware attribute selection to construct a compact attribute supervision mask for an image-only hierarchical classifier.

Fifth, when multiple attribute groups are introduced into a visual model, it is necessary to understand which components actually contribute to any improvement. Attribute-guided learning may improve performance because of auxiliary semantic supervision, because of species-level feature fusion, because of attribute-query cross-attention, or simply because of additional model capacity. Therefore, controlled ablation studies are required, including direct-fusion alternatives and variants without auxiliary attribute supervision. In addition, diversity-based regularization methods in fine-grained classification suggest that reducing redundancy among discriminative features can be useful \cite{Chang_2020}. In this thesis, this idea is adapted only as a regularization principle for attribute-group attention patterns; it is not used to claim pixel-level or part-level attribute localization.

Motivated by these gaps, this thesis investigates attribute-guided fine-grained hierarchical bird classification. The proposed approach uses CUB attributes as training-time auxiliary supervision, fuses attribute-oriented features into the species-level branch, compares attribute-query cross-attention with a direct-fusion baseline, and studies Hard-Sibling Discriminative Attribute Selection (HSDAS) as a compact attribute-selection strategy for distinguishing visually similar sibling species. Through comparisons with FullAttr-312, Sibling-IVAS-188, Random-199, and HSDAS variants, this thesis evaluates whether a compact and hard-sibling-aware attribute subset can better support species-level recognition while monitoring hierarchical metrics such as Full-Path Accuracy, weighted Average Precision, and Tree-based InConsistency Error.